\begin{textsample}{POS Dim 6 – vem\_ed\_18 – Score -2.00 – vem\_ed\_18/t084.txt}  \label{ex:f6_pos_081}
UM FAROL DE ESPERANÇA CRÍTICA Hope Windle, diretora do SUNY COIL Center (State University of New York / EUA), referência mundial em Intercâmbios Virtuais, encerrou o ciclo de \textbf{apresentações}.
Após a exibição de um vídeo com trechos da entrevista concedida por videoconferência para VEm 17 em 23 de março de 2023, Hope Windle respondeu a perguntas da audiência.
No vídeo gravado, Hope mencionou que mais pessoas do setor de design instrucional, de apoio pedagógico e da área de ensino têm vindo participar de projetos COIL na SUNY, enquanto no início da pandemia eram executivos internacionais e pessoas ligadas à internacionalização no currículo.
Sobre o que vislumbra para os próximos 5 anos, ela é entusiasmada com o que virá. “Nos anos 2000 a 2010, as pessoas que primeiro adotaram esses projetos se \textbf{perguntavam}: como podemos usar essas tecnologias?
Agora, as pessoas estão reconhecendo COIL, que não é algo para as pessoas sem condições econômicas que não podem viajar e sim como um farol de esperança crítica para que as pessoas possam repensar sua educação e o que vão fazer com suas vidas”.
Para contextualizar esse conceito em voga na contemporaneidade, fez a seguinte citação:"A esperança crítica reflete a habilidade de avaliar o ambiente de maneira realista, através de uma lente de equidade e justiça, ao mesmo tempo em que vislumbra a possibilidade de um futuro melhor"(DUGAN, 2017; DUNCAN- ANDRADE, 2009).
Hope mencionou ainda a questão dolorosa da escravidão que afeta discussões sobre classe e raça, tanto nos Estados Unidos quanto no Brasil.
Pontuou questões-chave para a educação global na atualidade e no futuro próximo, como o pensamento inclusivo e a interseccionalidade.
A diretora do SUNY COIL Center citou Abby Bryant, diretora do programa EOP, grupo de apoio acadêmico, psicológico e financeiro para pessoas com dificuldades socioeconômicas em universidades dos EUA.

% matched lemmas: apresentação, perguntar
\end{textsample}
